\newpage

\section{Forwarding toyRISC}

\subsection{Structure}

The top module is the same as for the pipelined toyRISC. The file {\tt 0\_DEFINES.vh} contains few additional instructions:
\begin{verbatim}
`define chsgn	6'b01_0000 // rf[d]<=~rf[l]+1;
`define max		6'b01_0001 // rf[d]<=max(rf[l],rf[r]);
`define min		6'b01_0010 // rf[d]<=min(rf[l],rf[r]);
`define asub	6'b01_0011 // rf[d]<=abs(rf[l]-rf[r]);
\end{verbatim}
and is missing one:
\begin{verbatim}
`define load   	6'b10_0111 // rf[d]<=dataOut;
\end{verbatim}
because the instruction {\tt read} includes the loading in the register file. The {\tt memorySystem} module is the same as for the pipelined toyRISC. 

\includeverilogcode[caption={Forwarded toyRISC processor. File: toyRISCforwarded.sv.  (SystemVerilog)}, 
                    label={lst:toyRISCforwarded}, 
                    firstline=1, 
                    firstnumber=1, 
                    lastline=450]
                    {\verilogSourceDir/A17_toyRISCforwarded.sv}

\includeverilogcode[caption={Decode for the forwarded toyRISC processor. File: 3\_DECODE.sv.  (SystemVerilog)}, 
                    label={lst:DECODEf}, 
                    firstline=1, 
                    firstnumber=1, 
                    lastline=450]
                    {\verilogSourceDir/A17_DECODEf.sv}

\includeverilogcode[caption={Interrupt automaton for the forwarded toyRISC processor. File: 4\_intAut.sv.  (SystemVerilog)}, 
                    label={lst:intAut}, 
                    firstline=1, 
                    firstnumber=1, 
                    lastline=450]
                    {\verilogSourceDir/A17_intAut.sv}

The first level of the pipeline, {\tt 3\_INSTRFETSH}, is identic with the first level in the previous version, the pipelined version without forwarding..

\includeverilogcode[caption={Operands fetch for the forwarded toyRISC processor. File: 3\_OPSFETCH.sv.  (SystemVerilog)}, 
                    label={lst:OPSFETCHf}, 
                    firstline=1, 
                    firstnumber=1, 
                    lastline=450]
                    {\verilogSourceDir/A17_OPSFETCHf.sv}

\includeverilogcode[caption={{\tt intUnit} for the forwarded toyRISC processor. File: 4\_intUnit.sv.  (SystemVerilog)}, 
                    label={lst:intUnitF}, 
                    firstline=1, 
                    firstnumber=1, 
                    lastline=450]
                    {\verilogSourceDir/A17_intUnitF.sv}

\includeverilogcode[caption={Execute for the forwarded toyRISC processor. File: 3\_EXECUTE.sv.  (SystemVerilog)}, 
                    label={lst:EXECUTEf}, 
                    firstline=1, 
                    firstnumber=1, 
                    lastline=450]
                    {\verilogSourceDir/A17_EXECUTEf.sv}

\includeverilogcode[caption={ALU for the forwarded toyRISC processor. File: 4\_ALUf.sv.  (SystemVerilog)}, 
                    label={lst:ALUf}, 
                    firstline=1, 
                    firstnumber=1, 
                    lastline=450]
                    {\verilogSourceDir/A17_ALUf.sv}

\includeverilogcode[caption={Data memory stage of pipe for the forwarded toyRISC processor. File: 3\_DATAMEM.sv.  (SystemVerilog)}, 
                    label={lst:DATAMEM}, 
                    firstline=1, 
                    firstnumber=1, 
                    lastline=450]
                    {\verilogSourceDir/A17_DATAMEM.sv}






\subsection{Code generator}

The task for the instruction {\tt LOAD} must be removed , and the following four task added:


\includeverilogcode[caption={}, 
                    label={}, 
                    firstline=1, 
                    firstnumber=1, 
                    lastline=450]
                    {\verilogSourceDir/A17_ADDEDTASK.sv}






\subsection{Simulator}


\includeverilogcode[caption={Code generator for the forwarded toyRISC processor. File: 0\_testForvardedRISC.sv.  (SystemVerilog)}, 
                    label={lst:testForvardedRISC}, 
                    firstline=1, 
                    firstnumber=1, 
                    lastline=450]
                    {\verilogSourceDir/A17_testForvardedRISC.sv}




\subsection{Testing}

For testing procedure see \ref{testProcedure}.

\includeverilogcode[caption={Programs for the forwarded toyRISC processor. File: 0\_program.sv.  (SystemVerilog)}, 
                    label={lst:programF}, 
                    firstline=1, 
                    firstnumber=1, 
                    lastline=450]
                    {\verilogSourceDir/A17_programF.sv}
